\documentclass[a4paper]{article}
\usepackage{amsfonts}
\usepackage{amsmath}
\usepackage{amssymb}
\usepackage{graphicx}     
\usepackage{cancel}

% Command definities van frequent gebruikte commandos
\newcommand{\N}{\mathbb{N}}
\newcommand{\Q}{\mathbb{Q}}
\newcommand{\R}{\mathbb{R}}
\newcommand{\C}{\mathbb{C}}
\newcommand{\Bgsin}{\operatorname{Bgsin}}
\newcommand{\Bgcos}{\operatorname{Bgcos}}
\newcommand{\Bgtan}{\operatorname{Bgtan}}
\newcommand{\cis}{\operatorname{cis}}
\newcommand{\Limit}{\lim\limits_}

\begin{document}
  
\noindent \large Auteur: Wout Vanderborght\\
\noindent \large Studierichting: Informatica\\
\noindent \large Oefeningenreeks 2E nr. 26.16\\

\medskip

\normalsize

Los de volgende logaritmische vergelijking op: \\

$ 2\left( \log_2(x-1) - \log_4(x^2+1) \right) = 1 - \log_4 25 $ \\

$ \Leftrightarrow \cancel{2}\left( \dfrac{\log_2 (4) \cdot \log_2 (x-1) - \log_2(x^2+1)}{\cancel{\log_2 4}} \right) = 1 - \dfrac{\log_2 5^2}{\log_2 4} $ \\

$ \Leftrightarrow 2 \cdot \log_2 (x-1) - \log_2(x^2+1) = 1 - \log_2 5 $ \\

$ \Leftrightarrow \log_2 \left(\dfrac{(x-1)^2}{x^2+1}\right) = \log_2 2 - \log_2 5 $ \\

$ \Leftrightarrow \log_2 \left(\dfrac{x^2-2x + 1}{x^2+1}\right) = \log_2 \left(\dfrac{2}{5}\right) $ \\

$ \Leftrightarrow \dfrac{x^2-2x + 1}{x^2+1} = \dfrac{2}{5} $ \\

$ \Leftrightarrow 5(x^2-2x + 1) = 2(x^2+1) $ \\

$ \Leftrightarrow 5x^2 - 10x + 5 = 2x^2 + 2 $ \\

$ \Leftrightarrow 3x^2 - 10x + 3 = 0 $ \\

$ \Leftrightarrow x(3x-1) - 3(3x-1) = 0 $ \\

$ \Leftrightarrow (x-3)(3x-1) = 0 $ \\

$ x_1 = 3 $ \\

$ \cancel{x_2 = \dfrac{1}{3}} $ want $ x > 1 $, anders is het onbepaald als je het in de opgave invult. 

\begin{thebibliography}{999}
%\bibitem{a} Referentie
\end{thebibliography}
\end{document} 
